% Chapter 1 -- How to Use This Book.
% Sources (Law 1): docs/pubs/originals/ibm-dos.md Sec 8.1 (entry anatomy),
% Sec 9 (notation); PUB-REF-0001 AN-01, NT-01..NT-17; pubs/screens/disk/
% MANIFEST (the practice diskette); pubs/screens/capture.sh (pinned clock
% 1992-03-02 09:30:00); os/milton/kmain.c (start-up order); os/milton/
% command.c (prompt $P$G, drive A only: "B:" is not a command;
% startup_datetime: the date and time questions without AUTOEXEC.BAT,
% initech-qskz; the startup screen boots the practice diskette without its
% AUTOEXEC.BAT, screens/SCREENS disk field noautoexec).
% GATES: test-boot test-shell test-pubs-screens
\Gates{test-boot, test-shell, test-pubs-screens}
\chapter{How to Use This Book}

\section{How the Commands Are Described}
\index{command entries}
Chapters 2, 3 and 4 describe the commands of InitechDOS, one command to an
entry.  Each entry begins on a new page.  Its title gives the name of the
command and, in parentheses, the word the name is made from: \cmd{MD}
(Make Directory).  If an entry runs on to a second page, the title is
repeated at the top of that page.

Every entry has the same parts, always in the same order:

\SideHead{Purpose:}One sentence that says what the command does.

\SideHead{Format:}The exact form of the command, in the notation described
at the beginning of Chapter~2.

\SideHead{Type:}Whether the command is \emph{internal} (built into the
command processor and always available) or \emph{external} (a program file
on disk).  Three asterisks are printed under the word that applies.
Configuration commands and batch subcommands have no Type line.

\SideHead{Remarks:}What each parameter does, followed by the details you
need to use the command: what it displays, the limits on its use, and the
messages it can produce.

\SideHead{Example:}One or more examples, each introduced by a sentence that
says what it does.

\section{The Examples}
\index{examples}
Every example in this book that shows the display was produced by typing
the command shown at an InitechDOS system and recording what appeared.
What you type is shown in \type{bold type}; what InitechDOS displays is in
\texttt{regular type}.  The command prompt is shown as part of each line
you type, exactly as it appears on the screen:
\begin{TypeIn}
A:\>
\end{TypeIn}
Each example begins with a freshly started system and a copy of the
InitechDOS practice diskette in drive~A.  The examples were recorded with
the system clock set to 9:30 a.m.\ on March~2, 1992, so the dates you see
on your screen will differ from those in the book.

\section{The Practice Diskette}
\index{practice diskette}
The practice diskette has the volume label \file{INITECH} and holds these
files and directories:
\par\medskip
\begin{tabular}{@{}l@{\quad}p{66mm}@{}}
{\imSideHeadFont Name} & {\imSideHeadFont Contents}\\[2pt]
\file{AUTOEXEC.BAT} & Two commands carried out when the system starts:
  \texttt{@ECHO OFF} and \texttt{SET PATH=A:\bs APPS}\\
\file{README.TXT}   & A note about the diskette\\
\file{MEMO.TXT}     & A short office memorandum\\
\file{PHONE.TXT}    & The internal telephone list\\
\file{APPS}         & A directory holding the utility programs
  \file{CHKDSK.COM}, \file{ATTRIB.COM}, \file{LABEL.COM}, \file{XCOPY.COM},
  \file{TREE.COM} and \file{EDLIN.COM}\\
\file{SAMPLES}      & A directory holding the batch files used in
  Chapter~4\\
\end{tabular}
\par\medskip
Because \file{AUTOEXEC.BAT} names the directory \file{\bs APPS} in the
\cmd{PATH} environment variable, you can run the utility programs from any
directory by typing their names.  Use a copy of the practice diskette for
the examples, not the original: several examples change the files on it.

\section{Starting InitechDOS}
\index{starting InitechDOS}
When the computer is turned on, InitechDOS displays its name and version
and the copyright notice, reads \file{CONFIG.SYS} (Chapter~3), carries out
\file{AUTOEXEC.BAT} (Chapter~4), and displays the command prompt.  When
there is no \file{AUTOEXEC.BAT} file, InitechDOS instead displays the
current date and asks for a new one, then does the same for the time,
before it displays the prompt.  Press \key{Enter} at each question to keep
the date or time shown, or type a new one as for the \cmd{DATE} and
\cmd{TIME} commands; an answer that is not a valid date or time displays
\texttt{Invalid date} or \texttt{Invalid time} and the question again.
Started from the practice diskette with its \file{AUTOEXEC.BAT} removed,
the screen looks like this:
\Shot{startup}

\section{Drives}
\index{drives}
In this version of InitechDOS your files are kept on one disk, drive~A,
whether it is a diskette or a \Slip{VO30-11}{hard disk}, and the command prompt always begins with \file{A:}.  Where a format in this
book shows a drive letter, \ph{d}\kw{:}, the drive is drive~A; you may
type \type{A:} or leave the drive out.  Typing another drive letter, such
as \type{B:}, at the prompt displays \texttt{Bad command or file name}.
